\documentclass[UTF8,12pt,a4paper]{ctexart}

% ==================== 基础宏包 ====================
\usepackage{geometry}
\geometry{left=2.5cm,right=2.5cm,top=2.5cm,bottom=2.5cm}
\usepackage{fancyhdr}       % 页眉页脚
\usepackage{tocloft}        % 专门用来定制目录样式（解决点点点问题）

% ==================== 定制目录的点线样式 ====================
% 强制给一级标题、二级标题、三级标题都加上虚线引导线
\renewcommand{\cftsecleader}{\cftdotfill{\cftdotsep}}
\renewcommand{\cftsubsecleader}{\cftdotfill{\cftdotsep}}
\renewcommand{\cftsubsubsecleader}{\cftdotfill{\cftdotsep}}

% ==================== 超链接设置（必须在最后导入） ====================
% hidelinks 参数可以去掉所有的红框和颜色，保留点击跳转功能
\usepackage[hidelinks]{hyperref} 

% ==================== 页眉页脚设置 ====================
\pagestyle{fancy}
\fancyhf{}
\cfoot{\thepage}
\setlength{\headheight}{15pt}

\begin{document}

% ==================== 1. 封面（无页码） ====================
\title{\textbf{基于深度学习的图像识别研究}}
\author{张三\\某某大学计算机学院}
\date{\today}
\maketitle
\thispagestyle{empty}       % 封面不显示页码

% ==================== 2. 前置部分（罗马数字页码） ====================
\newpage
\pagenumbering{roman}       % 开始用罗马数字 i, ii, iii...

% --- 中文摘要 ---
\section*{摘要}
\addcontentsline{toc}{section}{摘要} % 把“摘要”加入目录
这是中文摘要。研究生学位论文的摘要通常要写好几段。

\vspace{1em}
\noindent\textbf{关键词：} LaTeX；论文排版；深度学习

\newpage

% --- 英文摘要 ---
\section*{Abstract}
\addcontentsline{toc}{section}{Abstract}
This is the English abstract. You can write as many paragraphs as you need.

\vspace{1em}
\noindent\textbf{Keywords:} LaTeX; Thesis; Deep Learning

\newpage

% --- 目录 ---
\tableofcontents
\addcontentsline{toc}{section}{目录}
\newpage

% ==================== 3. 正文部分（阿拉伯数字页码） ====================
\pagenumbering{arabic}      % 正文从 1 开始用阿拉伯数字

\section{引言}
\label{sec:intro}
这是正文的第一章。

\subsection{研究背景}
这是二级标题。
\subsubsection{国内外现状}
这是三级标题。不要滥用，一般到三级就够了。

\section{相关技术}
\label{sec:related}
这是第二章。

\subsection{卷积神经网络}
这是二级标题。

\section{结论}
\label{sec:conclusion}
总结全文。

\end{document}